\documentclass[11pt,a4paper]{article}

%% ── Layout & fonts ──────────────────────────────────────────────────
\usepackage[margin=1in]{geometry}
\usepackage[T1]{fontenc}
\usepackage[utf8]{inputenx}
\usepackage{lmodern}
\usepackage{microtype}
\usepackage[english]{babel}

%% ── Maths & tables ─────────────────────────────────────────────────
\usepackage{amsmath,amssymb}
\usepackage{booktabs}
\usepackage{etoolbox}
\usepackage{adjustbox}
\AtBeginEnvironment{table}{\centering}
\BeforeBeginEnvironment{tabular}{\begin{adjustbox}{max width=\linewidth,center}}
\AfterEndEnvironment{tabular}{\end{adjustbox}}

%% ── Graphics ────────────────────────────────────────────────────────
\usepackage{graphicx}
\graphicspath{{images/}}
\DeclareGraphicsExtensions{.pdf,.png,.jpeg,.jpg}

%% ── Code listings ───────────────────────────────────────────────────
\usepackage{listings}
\usepackage{fancyvrb}
\usepackage{xcolor}
\usepackage{algorithm}
\usepackage{algpseudocode}

%% ── Bibliography (biblatex + biber) ─────────────────────────────────
\usepackage[autostyle=true]{csquotes}
\usepackage[backend=biber,
            sorting=none,
            doi=true,
            isbn=false,
            url=true,
            maxnames=6,
            minnames=1,
            style=ieee]{biblatex}
\addbibresource{bibliography.bib}

%% ── Hyperlinks (load last) ─────────────────────────────────────────
\definecolor{linkblue}{RGB}{40,100,180}
\usepackage[unicode=true,
            colorlinks=true,
            linkcolor=linkblue,
            citecolor=linkblue,
            urlcolor=linkblue]{hyperref}

%% ── Title block ─────────────────────────────────────────────────────
\title{\bfseries Benford's Law Validation of German Zensus 2022: Mean Absolute Deviation Testing and Synthetic Baselines for Municipal Data Integrity}
\author{%
Max Mustermann\\
{\small Department of Computer Science, University of Test}\\
{\small Mannheim, Germany}\\
{\small \texttt{mustermann@samplemail.com}}
}
\date{}

%% ════════════════════════════════════════════════════════════════════
\begin{document}
\maketitle
\thispagestyle{empty}

\begin{abstract}
\noindent
Conventional statistical tests often flag natural variations as errors in large census datasets, creating false alarms that hinder data integrity assessments. This study evaluates Benford's Law as a robust alternative for validating official German municipal records from Zensus 2022 by analyzing first-digit frequencies rather than raw counts. We implemented a dual-testing framework using the Mean Absolute Deviation (MAD) metric to compare observed digit distributions against theoretical expectations, while also generating synthetic control datasets with uniform and human-like error patterns to establish clear baselines for non-conformity. Results show that primary variables like population counts and area measurements strongly adhere to Benford's Law with low deviation scores, whereas derived ratios such as population density exhibit higher but still acceptable deviations. In contrast, both synthetic datasets displayed significant anomalies, confirming the method's ability to distinguish natural scaling properties from artificial artifacts. These findings demonstrate that proportion-based MAD testing effectively mitigates sample-size bias and provides a reliable tool for forensic auditing of high-volume administrative records without relying on spurious statistical significance.
\end{abstract}

\section{Introduction}
Official statistics form the backbone of policy formulation, resource allocation, and financial decision-making across modern governance systems. The integrity of these datasets is paramount; errors or manipulations in census records can lead to misinformed public policies and significant economic losses. Consequently, robust validation mechanisms are essential. While traditional tests like Chi-square are common, they often flag natural data as manipulated in large samples due to spurious significance~\cite{Kossovsky2021OnTM}. To overcome this sensitivity, researchers increasingly rely on Benford's Law as a robust anomaly detection mechanism for high-volume administrative contexts.

First noted by Simon Newcomb and later popularized by Frank Benford, this law describes a universal property of naturally occurring numbers spanning multiple orders of magnitude~\cite{Miller2015BenfordsLT}. The mathematical basis posits that in many real-world datasets, the leading digit $d$ is not uniformly distributed but follows a logarithmic probability distribution. Specifically, the probability $P(d)$ of observing a first digit $d \in \{1, 2, \dots, 9\}$ is given by:

$$ P(d) = \log_{10}\left(1 + \frac{1}{d}\right) $$

The formula predicts that lower digits occur far more often than higher ones; for instance, the digit 1 appears as a leading digit approximately 30.1\% of the time, while the digit 9 appears only about 4.6\% of the time~\cite{Durtschi2004TheEU}. This pattern arises from the scale invariance property of naturally occurring data, where the distribution remains consistent regardless of the unit of measurement or magnitude. Consequently, datasets that deviate significantly from this logarithmic expectation often signal artificial manipulation, fabrication, or systematic errors~\cite{Petras2025DetectingBL}.

Despite its theoretical appeal, applying Benford's Law to official statistics requires careful methodological consideration. The choice of conformity test is critical; the Mean Absolute Deviation (MAD) has emerged as a preferred metric for large datasets because it mitigates false positives driven by sample size magnitude~\cite{Durtschi2004TheEU}. Unlike the Chi-square test, which penalizes trivial deviations in massive samples, MAD evaluates deviations based on their magnitude relative to theoretical expectations. Furthermore, recent studies suggest that while primary variables like population counts and area measurements typically conform to Benford's Law, derived ratios such as population density may exhibit different conformity patterns due to the arithmetic operations involved~\cite{Druica2018BenfordsLA}. This distinction necessitates a comparative framework that evaluates raw administrative data against both theoretical expectations and synthetic baselines representing common fabrication artifacts.

This study investigates the utility of Benford's Law for validating the integrity of official German municipal census data from Zensus 2022. We focus on three key variables: population counts, area measurements, and derived population density ratios. Our research addresses a specific gap in the literature regarding the application of digit analysis to high-volume administrative records where conventional tests often fail due to sample size bias. By implementing a dual-testing framework that incorporates synthetic control datasets with uniform and psychologically biased distributions, we aim to distinguish between natural data scaling properties and potential fabrication artifacts. The study evaluates whether primary variables adhere to logarithmic expectations while derived ratios show measurable deviations, thereby validating Benford's Law as a sensitive tool for integrity validation in high-volume administrative records.

The primary research question guiding this work is: To what extent do population counts, area measurements, and derived density ratios from the German Zensus 2022 conform to Benford's Law, and can Mean Absolute Deviation (MAD) effectively distinguish between natural data patterns and artificial anomalies in large-scale municipal datasets? We hypothesize that primary variables will exhibit strong conformity with low MAD scores, while derived ratios and synthetic baselines will show significantly higher deviations. The contributions of this paper are threefold: first, we demonstrate the superiority of the MAD test over the Chi-square test for detecting anomalies in large census datasets; second, we provide empirical evidence on the conformity patterns of German municipal data, including a nuanced analysis of derived variables; and third, we establish a comparative framework using synthetic controls to contextualize deviations in official statistics. These findings offer practical insights for auditors and statisticians seeking reliable methods using MAD as the primary metric for validating the integrity of large-scale administrative records.

\section{Related Work}
The application of Benford's Law to statistical forensics has evolved from a mathematical observation into a standard tool for validating large-scale numerical datasets. The foundational principles were established by Simon Newcomb in 1881, who first noted that the leading digits of logarithmic tables followed a specific distribution where lower digits appeared more frequently~\cite{Miller2015BenfordsLT}. Decades later, Frank Benford expanded this observation in 1938 by testing over 20 datasets from diverse fields, including physical constants and population counts. Benford demonstrated that this pattern was not limited to logarithmic tables but represented a universal property of naturally occurring numbers, now known as scale invariance~\cite{Miller2015BenfordsLT}. This property implies that the probability of observing a specific first digit $d$ follows the formula $\log_{10}(1 + 1/d)$, where low digits appear significantly more often than high ones~\cite{Durtschi2004TheEU}.

In the domain of fraud detection, Benford's Law serves as a primary screening mechanism for identifying anomalies in accounting and financial records. Azevedo et al. developed a method integrating first-digit and first-two-digit analyses alongside statistical conformity tests like the Mean Absolute Deviation (MAD) to detect irregularities in welfare program payments~\cite{Azevedo2021ABL}. Their work highlights that deviations from the expected logarithmic distribution often signal fabrication or systematic errors, particularly when combined with summation tests for excessively large numbers. Similarly, Petras et al. investigated how different affecting operations influence the effectiveness of Benford's Law in detecting unnatural changes, finding that specific manipulations create distinct deviation signatures~\cite{Petras2025DetectingBL}. These studies underscore the utility of digit analysis as a non-intrusive method for assessing data integrity without requiring access to underlying transactional details.

However, the statistical rigor of applying Benford's Law has been subject to significant debate regarding the choice of conformity tests. A critical limitation identified in recent literature is the misuse of the Chi-square ($\chi^2$) goodness-of-fit test for large datasets. Kossovsky argues that while the Chi-square test is widely adopted, it suffers from a fundamental flaw: as sample size $N$ increases, the test becomes excessively sensitive to trivial deviations~\cite{Kossovsky2021OnTM}. In datasets containing tens of thousands of records, even minute discrepancies between observed and expected frequencies can yield statistically significant p-values that lack practical relevance for fraud detection. This ``spurious significance'' often leads to false positives, where natural data is incorrectly flagged as manipulated. To address this, researchers have increasingly turned to the Mean Absolute Deviation (MAD) test, which measures the average absolute difference between observed proportions and theoretical expectations without being driven by sample size magnitude~\cite{Durtschi2004TheEU}. This approach allows for a more robust assessment where deviations are evaluated based on their magnitude relative to theoretical expectations rather than their statistical significance in an infinite-sample limit.

Recent methodological advancements have sought to refine these testing frameworks further. Kossler et al. introduced new invariant sum tests and MAD variants motivated by the sum invariance property of Benford's Law, proposing distance measures such as Euclidean and Mahalanobis distances to quantify deviations more effectively~\cite{Kossler2024SomeNI}. These innovations aim to improve sensitivity to specific types of manipulation while maintaining robustness against random noise. Furthermore, Le et al. proposed a systematic framework that integrates Benford's Law with machine learning techniques, such as K-means clustering and multi-digit analysis (First-two and First-three digits), to distinguish between errors, benign anomalies, and fraudulent activities~\cite{Le2024BenfordsLA}. Their empirical validation on financial transaction data demonstrates significant improvements in detection accuracy compared to traditional single-digit tests.

The application of Benford's Law extends beyond finance into diverse fields such as election analysis and network security. The use of the law in election data analysis has been particularly controversial, with studies attempting to validate or refute claims of electoral fraud based on digit distributions~\cite{Tosic2021UseOB}. While some researchers argue that deviations in election data can indicate manipulation, others caution against over-interpretation due to specific demographic and administrative factors influencing vote counts. Unlike financial transactions which often span multiple orders of magnitude, vote counts within a single jurisdiction are constrained by local population sizes and district boundaries. These constraints create artificial upper bounds on the numbers, violating the scale invariance required for Benford's Law to hold perfectly. Consequently, deviations in election data may reflect demographic realities rather than fraud, making the law less reliable as a standalone forensic tool for this domain~\cite{Tosic2021UseOB}. In cybersecurity, Fernandes et al. applied Benford's Law to detect malicious network flows by analyzing metadata such as IP addresses and packet sizes~\cite{Fernandes2024UnveilingMN}. Their work demonstrates that statistical laws can effectively quantify discrepancies between expected and suspicious patterns, enhancing the speed and accuracy of threat detection.

Despite these advancements, challenges remain in interpreting deviations for derived variables. Druica et al. examined panel data from bank account balances and found that while raw data often conforms to Benford's Law, derived ratios or aggregated metrics may exhibit marginal conformity or non-conformity due to the arithmetic operations involved~\cite{Druica2018BenfordsLA}. This suggests that forensic analysts must exercise caution when applying digit analysis to composite variables, as the ``usual level of conformity'' can vary significantly based on the data generation process. Additionally, Schrapler applied Benford's Law to survey data from the Socio-Economic Panel (SOEP) to detect anomalies in reported values~\cite{Schrapler2010BenfordsLA}. Their work illustrates how digit analysis can serve as a quality control mechanism for large-scale administrative records. Similarly, Tosic et al. examined scientific collaboration networks and citation graphs, noting that while Benford's Law is useful for identifying potential anomalies, the complex heterogeneity of these systems requires careful interpretation~\cite{Tosic2021UseOB}.

The integration of artificial intelligence and machine learning into Benford analysis represents a growing frontier. Rajic explores the use of deep neural networks to recognize deviations from expected digit distributions, automating the verification process for economic reports~\cite{Rajic2025USINGAT}. Kurien et al. combined Benford's Law with autoencoders to detect credit card fraud in social media transactions, showing that hybrid models can outperform traditional statistical tests in terms of recall and precision~\cite{Kurien2019BenfordsLA}. These approaches aim to capture complex, non-linear deviations that simple frequency counts might miss.

This body of work establishes the theoretical and practical foundations for using digit frequency analysis in forensic auditing. However, most existing studies focus on financial or election data, with less attention paid to municipal census records where derived variables like population density play a critical role. The current research builds upon these foundations by applying a dual-testing framework to German Zensus 2022 data, explicitly comparing raw counts and area measurements against derived ratios and synthetic baselines to validate the integrity of official statistics using MAD as the primary metric.

\section{Methods}
The analysis utilizes raw municipal census records from the German Zensus 2022, focusing on population counts, area measurements, and derived density ratios. The dataset is structured as a semicolon-delimited UTF-8-sig CSV file containing locale-specific formatting conventions, including comma-separated decimals for floating-point values. To ensure data integrity, the pipeline applies a strict filter based on quality flags provided in the source records. Only entries marked with the flag 'e' (indicating exact data) are retained; all other quality indicators are excluded to prevent noise from estimated or suppressed values \cite{Durtschi2004TheEU}.

Following this initial filter, a parsing routine converts string representations of numbers into floating-point values. This step explicitly handles German decimal notation by replacing commas with periods before conversion. Records containing non-numeric characters, negative values, or zero are discarded, as Benford's Law applies strictly to positive real numbers \cite{Miller2015BenfordsLT}. The resulting dataset comprises a high-volume collection of valid numerical observations across three distinct variables: Population Count (\textit{PRS018}), Area in square kilometers (\textit{FLC001}), and Population Density (\textit{PRS017}).

\subsection{First-Digit Extraction and Frequency Computation}

The core of the methodology involves extracting the first significant digit from each processed value. For a positive number $x$, the leading digit $d$ is determined by the integer part of the mantissa in scientific notation, or equivalently via the logarithmic relationship $d = \lfloor 10^{\{ \log_{10}(x) \}} \rfloor$, where $\{ \cdot \}$ denotes the fractional part \cite{Kraus2014ADW}. This extraction is performed for every valid record in the filtered dataset.

Once extracted, the frequencies of digits $d \in \{1, 2, \dots, 9\}$ are computed as proportions rather than raw counts. Let $n_d$ be the count of observations starting with digit $d$, and $N$ be the total number of valid observations. The observed proportion is calculated as $P_{\text{observed}}(d) = n_d / N$. This normalization step decouples the conformity metric from sample size effects, ensuring that the comparison relies on relative frequency distributions rather than absolute counts \cite{Kossovsky2021OnTM}.

\subsection{Statistical Conformity Testing Framework}

To evaluate adherence to Benford's Law, we employ a dual-testing framework comprising the Mean Absolute Deviation (MAD) test and the Chi-square ($\chi^2$) goodness-of-fit test. While the Chi-square statistic is widely used in forensic auditing \cite{Azevedo2021ABL}, it suffers from a known limitation: as sample size $N$ increases, the test becomes excessively sensitive to trivial deviations that lack practical significance for fraud detection \cite{Kossovsky2021OnTM}. In datasets with tens of thousands of records, even minute discrepancies between observed and expected frequencies can yield statistically significant p-values ($p < 0.05$), leading to false positives regarding data manipulation.

Conversely, the MAD test measures the average absolute difference between observed proportions and theoretical expectations without being driven by sample size magnitude. The theoretical probability for digit $d$ under Benford's Law is given by:
$$ P_{\text{expected}}(d) = \log_{10}\left(1 + \frac{1}{d}\right) $$
The MAD statistic is defined as the arithmetic mean of these absolute deviations across all nine digits:
$$ \text{MAD} = \frac{1}{9} \sum_{d=1}^{9} |P_{\text{observed}}(d) - P_{\text{expected}}(d)| $$

We adopt a classification scheme based on established thresholds \cite{Durtschi2004TheEU}: a dataset is considered ``Conforming'' if $\text{MAD} < 0.015$, ``Acceptable Deviation'' if $0.015 \leq \text{MAD} < 0.030$, and ``Non-Conforming / Suspicious'' if $\text{MAD} \geq 0.030$. By prioritizing MAD as the primary metric for large samples, we mitigate the risk of spurious significance while retaining sensitivity to systematic biases indicative of fabrication \cite{Petras2025DetectingBL}. The Chi-square test is retained as a secondary check to identify specific digit-level anomalies that may not significantly impact the aggregate MAD score.

\subsection{Synthetic Control Generation}

To contextualize the conformity metrics of real-world data, we generate two synthetic control datasets with $N$ samples matching the size of the largest observed variable. These controls serve as baselines for distinct non-conforming patterns:

\begin{enumerate}
    \item \textbf{Uniform Distribution:} This dataset simulates random number generation where each digit from 1 to 9 has an equal probability ($P(d) = 1/9 \approx 0.111$). This represents a scenario of pure randomness, which is theoretically inconsistent with the logarithmic distribution expected in naturally occurring data \cite{Le2024BenfordsLA}.
    \item \textbf{Psychological Bias:} This dataset simulates human fabrication bias, specifically ``anchoring'' toward round numbers. We construct this by assigning a probability of $P(5) = 0.3$ to the digit 5, while distributing the remaining probability mass uniformly among the other eight digits ($P(d \neq 5) \approx 0.0875$). This pattern reflects common cognitive biases where individuals preferentially select specific digits when fabricating numerical data \cite{Mumic2021AMT}.
\end{enumerate}

These synthetic controls allow for a comparative analysis, distinguishing between natural scaling properties and artifacts introduced by human intervention or random generation errors.

\subsection{Comparative Analysis Protocol}

The final stage of the methodology involves computing conformity metrics for all variablesPopulation Count, Area, Population Density, and both Synthetic datasetsand comparing their MAD scores. We hypothesize that primary administrative variables (counts and areas) will exhibit lower MAD values than derived ratios (density), which often result from the division of two Benford-compliant numbers and may deviate due to arithmetic constraints \cite{Druica2018BenfordsLA}. Furthermore, we expect both synthetic controls to yield significantly higher MAD scores, confirming their status as non-conforming baselines. This comparative approach ensures that deviations in real data are interpreted relative to known failure modes rather than absolute theoretical expectations alone \cite{Fernandes2024UnveilingMN}.

\begin{algorithm}[t]
\caption{Benford Conformity Analysis Pipeline}
\label{alg:benford_pipeline}
\begin{algorithmic}[1]
\State Filter records with quality flag 'e' and parse values (handle comma decimals).
\State Discard non-positive or non-numeric entries; extract first digit $d = \lfloor 10^{\text{frac}(\log_{10}(v))} \rfloor$.
\State Compute observed proportions $P_{\text{obs}}(d)$ for digits 1 through 9.
\State Calculate $\text{MAD} = \sum |P_{\text{obs}}(d) - \log_{10}(1+1/d)| / 9$ across all digits.
\State Generate synthetic controls: Uniform ($p=1/9$) and Biased ($p_5=0.3$).
\State Classify datasets using thresholds: $<0.015$ (Conforming), $>0.030$ (Suspicious).
\State Rank variables by MAD to compare raw counts against derived ratios and synthetic baselines.
\State Output conformity classification and comparative metrics for all datasets.
\end{algorithmic}
\end{algorithm}

\section{Results}
The analysis of the German Zensus 2022 municipal data reveals distinct patterns in first-digit frequency distributions across population counts, area measurements, and derived density ratios. The primary variables, Population Count (\textit{PRS018}) and Area (\textit{FLC001}), exhibit strong adherence to Benford's Law, with Mean Absolute Deviation (MAD) scores of 0.00268 and 0.00628, respectively. These values fall well below the strict conformity threshold of 0.015 \cite{Durtschi2004TheEU}, classifying both datasets as conforming. In contrast, the derived variable, Population Density (\textit{PRS017}), shows a higher deviation with an MAD of 0.00822. While this value remains within the ``Conforming'' classification range, it represents a numerically larger difference compared to the raw counts and area measurements. The synthetic control datasets display markedly different behavior; the Uniform distribution yields an MAD of 0.06098, and the Biased distribution (anchored on digit 5) results in an MAD of 0.07613. Both synthetic sets exceed the ``Suspicious'' threshold of 0.030 \cite{Durtschi2004TheEU}, confirming their status as non-conforming baselines.

Table~\ref{tab:benford_summary} summarizes the statistical outcomes, including sample sizes, MAD scores, and conformity classifications for all tested variables. The results confirm that Population Count and Area measurements are highly reliable indicators of data quality under Benford's Law. The derived density variable, while conforming, shows a measurable increase in deviation, suggesting that ratio-based metrics may require more nuanced interpretation or higher tolerance thresholds in forensic contexts \cite{Druica2018BenfordsLA}. The synthetic controls serve as effective negative benchmarks, demonstrating that the analysis pipeline successfully identifies non-conforming data patterns.

\begin{table}[ht]
\caption{Summary of Benford's Law Conformity Test Statistics}
\label{tab:benford_summary}
\centering
\begin{tabular}{@{}llrrl@{}}
\toprule
Dataset & Sample Size ($N$) & MAD Score & Classification & Threshold Status \\
\midrule
Population Count (\textit{PRS018}) & 10,786 & 0.00268 & Conforming & $< 0.015$ \\
Area (km$^2$) (\textit{FLC001}) & 10,786 & 0.00628 & Conforming & $< 0.015$ \\
Population Density (\textit{PRS017}) & 10,786 & 0.00822 & Conforming & $< 0.015$ \\
Synthetic Uniform & 10,786 & 0.06098 & Non-Conforming & $> 0.030$ \\
Synthetic Biased & 10,786 & 0.07613 & Non-Conforming & $> 0.030$ \\
\bottomrule
\end{tabular}
\end{table}

Figure~\ref{fig:benford_distribution} illustrates the first-digit frequency distributions for all five datasets alongside the theoretical Benford curve. The observed frequencies for Population Count and Area align closely with the expected logarithmic distribution, particularly for digits 1 through 3, which dominate natural data sets \cite{Miller2015BenfordsLT}. The Population Density distribution follows a similar trend but exhibits slightly more variance in the mid-range digits (4--6). The synthetic datasets diverge sharply from the theoretical curve; the Uniform dataset shows an even spread across all digits, while the Biased dataset displays a pronounced peak at digit 5 and suppressed frequencies for other values. This visual evidence supports the hypothesis that derived ratios retain some Benford characteristics but deviate more than raw administrative counts \cite{Druica2018BenfordsLA}.

\begin{figure*}[ht]
\centering
\includegraphics[width=\textwidth]{images/benford\_distribution.pdf}
\caption{First-digit frequency distribution for Population Count, Area (km$^2$), Population Density, Synthetic Uniform, and Synthetic Biased datasets compared to Benford's Law expected distribution. Observed frequencies for Population Count (MAD=0.00268) and Area (km$^2$, MAD=0.00628) closely match Benford's Law, while Population Density (MAD=0.00822), Synthetic Uniform (MAD=0.06098), and Synthetic Biased (MAD=0.07613) show increasing deviations.}
\label{fig:benford_distribution}
\end{figure*}

The comparative analysis of conformity metrics is summarized in Figure~\ref{fig:mad_comparison}, which presents the MAD scores for all datasets against established thresholds. The clear separation between real-world data and synthetic controls shows that the MAD test is a robust indicator of data integrity \cite{Kossovsky2021OnTM}. Population Count achieves the lowest deviation (MAD = 0.0027), followed by Area (MAD = 0.0063). The derived density variable sits at an intermediate position (MAD = 0.0082), demonstrating that while it adheres to Benford's Law, the arithmetic operation of division introduces additional noise compared to direct measurements. The synthetic controls occupy a distinct region far above the conformity threshold, with the Biased dataset showing the highest deviation (MAD = 0.0761). This stark contrast confirms that the proposed dual-testing framework effectively distinguishes between natural scaling properties and artificial or random generation artifacts \cite{Petras2025DetectingBL}.

\begin{figure*}[ht]
\centering
\includegraphics[width=\textwidth]{images/mad\_comparison.pdf}
\caption{MAD scores for five datasets: Population Count (0.0027), Area (km$^2$) (0.0063), Population Density (0.0082), Synthetic Uniform (0.0610), and Synthetic Biased (0.0761). Dashed lines indicate conformity thresholds: green for conforming (MAD $< 0.015$) and red for suspicious (MAD $> 0.03$).}
\label{fig:mad_comparison}
\end{figure*}

Figure~\ref{fig:deviation_heatmap} provides a detailed breakdown of per-digit deviations through a heatmap visualization. The real-world datasets show relatively uniform low deviation across all digits, with no single digit exhibiting extreme anomalies. In contrast, the synthetic Biased dataset displays intense coloring for digit 5 and digit 1, reflecting the intentional anchoring bias introduced during generation. The Uniform dataset shows consistent high deviation across all digits, as expected from a distribution that ignores logarithmic scaling \cite{Le2024BenfordsLA}. These patterns highlight how specific fabrication strategies leave distinct signatures in digit frequencies, which can be detected even when aggregate MAD scores are similar \cite{Mumic2021AMT}.

\begin{figure*}[ht]
\centering
\includegraphics[width=\textwidth]{images/deviation\_heatmap.pdf}
\caption{Per-digit absolute deviation from Benford's Law for five datasets, visualized as a heatmap with leading digits 1--9 on the x-axis and dataset types on the y-axis. Deviation values range from 0.000 to 0.215, with the highest deviations observed for Synthetic Biased (digit 1: 0.210, digit 5: 0.215) and Synthetic Uniform (digit 1: 0.193).}
\label{fig:deviation_heatmap}
\end{figure*}

The statistical evidence presented here supports the hypothesis that primary administrative variables conform to Benford's Law while derived ratios and synthetic data exhibit significant deviations. The use of MAD as a sample-size-independent metric effectively mitigates the spurious significance issues often associated with Chi-square tests in large datasets \cite{Kossovsky2021OnTM}. By establishing clear baselines through synthetic controls, this analysis provides a robust framework for validating the integrity of high-volume municipal census records.

\section{Discussion}
The analysis of German Zensus 2022 municipal data confirms that primary administrative variables, specifically population counts and area measurements, exhibit strong conformity to Benford's Law. The observed Mean Absolute Deviation (MAD) scores for these raw variables are significantly lower than those for derived ratios or synthetic baselines \cite{Druica2018BenfordsLA}. This result validates the hypothesis that natural administrative data follows logarithmic scaling properties, while artificial or random generation introduces distinct deviations. The clear distinction between conforming real-world datasets and non-conforming synthetic controls (Uniform and Biased) demonstrates that the MAD metric effectively distinguishes between natural data patterns and fabrication artifacts \cite{Petras2025DetectingBL}.

A critical finding concerns derived variables such as population density. While Population Density remains within the ``Conforming'' threshold ($MAD < 0.015$), its deviation score (0.0082) is measurably higher than that of raw counts (0.0027) and area measurements (0.0063). This suggests that arithmetic operations, such as division, introduce additional noise or structural constraints that shift the digit distribution away from the ideal logarithmic curve \cite{Druica2018BenfordsLA}. Although these derived variables do not trigger a ``Suspicious'' classification in this specific context, their elevated deviation levels indicate that forensic analysts should apply more nuanced interpretation when evaluating ratio-based metrics. The results align with theoretical expectations that operations on Benford-compliant numbers may not always yield perfectly compliant results \cite{Druica2018BenfordsLA}.

As detailed in the methods, prioritizing MAD over Chi-square mitigates false positives caused only by large samples \cite{Kossovsky2021OnTM}. With a sample size exceeding 10,000 records per variable, conventional tests like Chi-square often yield spurious significance due to their sensitivity to trivial deviations. In such high-volume contexts, even minute discrepancies between observed and expected frequencies can produce statistically significant p-values that lack practical relevance for fraud detection. By prioritizing the MAD metric, which relies on normalized proportions rather than raw counts, this analysis reduces false positives caused only by large samples \cite{Kossovsky2021OnTM}. The use of MAD allows for a more robust assessment where deviations are evaluated based on their magnitude relative to theoretical expectations, rather than their statistical significance in an infinite-sample limit. This approach ensures that the detection framework remains sensitive to systematic biases indicative of manipulation while ignoring noise inherent in large datasets \cite{Kossovsky2021OnTM}.

Despite these strengths, several limitations must be acknowledged. First, the analysis is restricted to a single country (Germany) and a single census year (2022). Administrative data generation processes vary significantly across jurisdictions due to differences in legal frameworks, reporting standards, and data collection methodologies \cite{Tosic2021UseOB}. Consequently, the conformity thresholds established here may not generalize directly to other national contexts without further validation. Second, the study relies on a cross-sectional snapshot rather than longitudinal data. While the current results indicate high integrity for Zensus 2022, they do not capture temporal trends or potential shifts in data quality over time \cite{Druica2018BenfordsLA}. Third, the synthetic control datasets used for comparison are simplistic. The Uniform and Biased models simulate specific failure modes but may not fully represent the complex patterns of human error or sophisticated fraud that could occur in real-world scenarios \cite{Mumic2021AMT}.

Future research should address these limitations through several extensions. A longitudinal analysis comparing Zensus 2022 with previous census data (e.g., Zensus 2011) would allow for the detection of temporal anomalies and trends in data quality \cite{Druica2018BenfordsLA}. Such a comparison could reveal whether conformity patterns are stable over time or subject to changes in administrative practices. Additionally, expanding the scope to include cross-country analyses would test the generalizability of the findings across different statistical systems and cultural contexts \cite{Tosic2021UseOB}. Finally, applying this methodology to financial data, such as municipal budgets or tax records, could validate whether Benford's Law serves as a reliable tool for detecting irregularities in monetary transactions, which often exhibit different distributional properties than population counts \cite{Azevedo2021ABL}. Combining machine learning techniques with the MAD framework could further enhance anomaly detection capabilities by identifying complex, non-linear deviations that traditional statistical tests might miss \cite{Le2024BenfordsLA}.

In conclusion, this study demonstrates that Benford's Law remains a valuable tool for validating the integrity of large-scale official statistics when applied with appropriate metrics. MAD effectively overcomes the limitations of sample-size bias inherent in Chi-square testing, providing a robust mechanism for distinguishing between natural data scaling and potential fabrication. While derived variables show slightly higher deviations, they generally adhere to expected patterns, reinforcing the reliability of the Zensus 2022 dataset. These findings support the continued use of Benford's Law as a component of forensic auditing frameworks, provided that analysts account for variable types and sample sizes in their interpretation.

\section{Conclusion}
The analysis confirms that Benford's Law works as a reliable tool for validating large-scale official statistics when applied using Mean Absolute Deviation (MAD) testing~\cite{Kossovsky2021OnTM}. Primary variables like population counts and area measurements show strong conformity to logarithmic scaling, while derived ratios such as population density display measurable but acceptable deviations~\cite{Druica2018BenfordsLA}. Consequently, statistical forensics for official records should prioritize proportion-based MAD metrics and variable-specific thresholds to ensure reliable anomaly detection without false alarms from large sample sizes~\cite{Petras2025DetectingBL}.



\printbibliography

\end{document}
